% !TeX root = ../../infdesc.tex
\section{Variables and quantifiers}
\secbegin{secVariablesQuantifiers}


\begin{exercise}
\label{exEnglishToLogicalFormulae}
Find logical formulae that represent each of the following English statements.
\begin{enumerate}[(a)]
\item There is an integer that is divisible by every integer.
\item There is no greatest odd integer.
\item Between any two distinct rational numbers is a third distinct rational number.
\item If an integer has a rational square root, then that root is an integer.
\end{enumerate}
\end{exercise}

\begin{solutions*}
    \fixlistskip
    \begin{enumerate}[(a)]
\item $\exists q \in \mathbb{Z},\, \forall a \in \mathbb{Z},\, \exists k \in \mathbb{Z}, (q = ak)$
\item Define odd integer set $O=\{x \mid x=2k+1, k \in \mathbb{N}\}$
\[\neg \exists a \in O\, \forall b \in O, (a \ge b)\]
\item $\forall p,q \in \mathbb{R} \wedge p < q\, \exists r \in \mathbb{R}, (p < r \wedge r < q)$
\item $\forall a \in \mathbb{Z},\, \forall q \in \mathbb{Q},\, (q^2=a \implies q \in \mathbb{Z})$
\end{enumerate}
\end{solutions*}

\begin{exercise}
\label{exLogicalFormulaeToEnglish}
Find statements in plain English, involving as few variables as possible, that are represented by each of the following logical formulae. (The domains of discourse of the free variables are indicated in each case.)
\begin{enumerate}[(a)]
\item $\exists q \in \mathbb{Z},\, a = qb$ --- free variables $a, b \in \mathbb{Z}$
\item $\exists a \in \mathbb{Z},\, \exists b \in \mathbb{Z},\, (b \ne 0 \wedge bx = a)$ --- free variable $x \in \mathbb{R}$
\item $\forall d \in \mathbb{N},\, [(\exists q \in \mathbb{Z},\, n=qd) \Rightarrow (d = 1 \vee d = n)]$ --- free variable $n \in \mathbb{N}$
\item $\forall a \in \mathbb{R},\, [a > 0 \Rightarrow \exists b \in \mathbb{R},\, (b > 0 \wedge a < b)]$ --- no free variables
\end{enumerate}
\end{exercise}

\begin{solutions*}
    \fixlistskip
    \begin{enumerate}[(a)]
    \item $b$ divides $a$.
    \item $x$ is a rational number.
    \item The only positive divisors of $n$ are $1$ and $n$.
    \item There is no largest positive real number.
    \end{enumerate}
\end{solutions*}

\begin{exercise}
\label{exEveryIntegerIsRational}
Prove that every integer is rational.
\end{exercise}

\begin{solutions*}
    \fixlistskip
    \begin{proof}
        For any arbitrary integer $n$, let $p = n$ and $q = 1$, we get $p = qn$.

        Therefore, every integer is rational.
    \end{proof}
\end{solutions*}

\begin{exercise}
Prove that every linear polynomial over $\mathbb{Q}$ has a rational root.
\end{exercise}

\begin{solutions*}
    \fixlistskip
    \begin{proof}
        For any arbitrary linear polynomial $ax + b (a \neq 0)$ over $\mathbb{Q}$, it has a root $x = -\frac{b}{a}$. Since $a$ and $b$ are rational numbers, let $a = \dfrac{p}{q}$ and $b = \dfrac{r}{s}$, where $p, q, r, s \in \mathbb{Z}$ and $q, s \ne 0$. Then $x = -\dfrac{b}{a}=-\dfrac{rq}{sp}$ is a rational number.

        Therefore, every linear polynomial over $\mathbb{Q}$ has a rational root.
    \end{proof}
\end{solutions*}

\begin{exercise}
Prove that, for all real numbers $x$ and $y$, if $x$ is irrational, then $x+y$ and $x-y$ are not both rational.
\end{exercise}

\begin{solutions*}
    \fixlistskip
    \begin{proof}
    Suppose $x$ is irrational. Assume for the sake of contradiction that $x+y$ and $x-y$ are both rational.
    
    Let $a = x+y \in \mathbb{Q}$ and $b = x-y \in \mathbb{Q}$, then
    \[a+b = (x+y) + (x-y) = 2x\]
    \[x = \frac{a+b}{2} \in \mathbb{Q}\]
    This contracts the assumption that $x$ is irrational. 
    
    Therefore, $x+y$ and $x-y$ are not both rational.
    \end{proof}
\end{solutions*}

\begin{exercise}
Prove that there is a real number which is irrational but whose square is rational.
\end{exercise}

\begin{solutions*}
    \fixlistskip
    \begin{proof}
        We can prove this existentially quantified statements by constructing a specific example.

        Let $x = \sqrt{2}$, then $x$ is irrational and $x^2 = 2$ is rational.

        Therefore, there is a real number which is irrational but whose square is rational.
    \end{proof}
\end{solutions*}

\begin{exercise}
Prove that there is an integer which is divisible by zero.
\end{exercise}

\begin{solutions*}
    \fixlistskip
    \begin{proof}
        According to the Definition 0.33, $a$ is divisible $b$ if $a = qb$ for some integer $q$.

        Let $a = 0$ and $b = 0$, then $a = qb = 0q = 0$.

        Therefore, there is an integer which is divisible by zero.
    \end{proof}
\end{solutions*}

\begin{exercise}
\label{exExamplesOfUniqueExistentialQuantifier}
For each of the propositions, write it out as a logical formula involving the $\exists !$ quantifier and then prove it, using the structure of the logical formula as a guide.
\begin{enumerate}[(a)]
\item For each real number $a$, the equation $x^2+2ax+a^2=0$ has exactly one real solution $x$.
\item There is a unique real number $a$ for which the equation $x^2+a^2=0$ has a real solution $x$.
\item There is a unique natural number with exactly one positive divisor.
\end{enumerate}
\end{exercise}

\begin{solutions*}
    \fixlistskip
    \begin{enumerate}[(a)]
    \item $\forall a \in \mathbb{R},\, \exists! x \in \mathbb{R},\, x^2 + 2ax + a^2 = 0$
    \begin{proof}
        Let $a$ be an arbitrary real number. We must show that there exists a unique real number $x$ that satisfies $x^2 + 2ax + a^2 = 0$. We can rewrite the equation by factoring the perfect square:
        \[x^2 + 2ax + a^2 = (x+a)^2 = 0\]
        \begin{itemize}
            \item (Existence) Let $x = -a$. Since $a \in \mathbb{R}$, its negation $-a$ is also in $\mathbb{R}$. Substituting this into our factored equation yields:
            \[(-a + a)^2 = 0^2 = 0\]
            Thus, $x = -a$ is a valid real solution. This proves existence.
            \item (Uniqueness) Suppose there are two real solutions, $x_1$ and $x_2$. By our equation, this means:
            \[\begin{cases}
                (x_1 + a)^2 = 0 \implies x_1 = -a\\
                (x_2 + a)^2 = 0 \implies x_2 = -a
            \end{cases}\]
            Therefore, $x_1 = x_2$. This proves that the solution is unique.
        \end{itemize}
    \end{proof}
    \item $\exists! a \in \mathbb{R}, \exists x \in \mathbb{R}, x^2 + a^2 = 0$
    \begin{proof}
        We must show there is exactly one real number $a$ that allows the equation $x^2 + a^2 = 0$ to be solvable for a real number $x$.
        \begin{itemize}
            \item (Existence) Let $a = 0$. The equation becomes $x^2 + 0^2 = 0$, it has a real solutions $x = 0$. Thus, there exists at least one such real number $a$.
            \item (Uniqueness) Suppose there is some real number $b$ for which the equation $x^2 + b^2 = 0$ has a real solution $x'$. Because $x'$ and $b$ are real numbers, their squares must be non-negative:
            \[x'^2 \ge 0 \quad \text{and} \quad b^2 \ge 0\]
            The sum of two non-negative real numbers can only equal $0$ if both numbers are individually equal to $0$. Therefore, we must have:
            \[x'^2 = 0 \quad \text{and} \quad b^2 = 0\]
            Therefore, $b = a = 0$. This proves that the real number $a$ is unique.
        \end{itemize}
    \end{proof}
    \item $\exists! n \in \mathbb{N},\, (\exists! d \in \mathbb{N},\, \exists q \in \mathbb{N},\, n = qd)$
    \begin{proof}
        We must show there is exactly one natural number $n$ that has exactly one natural number $d$ dividing it.
        \begin{itemize}
            \item (Existence) Let $n = 1$. The positive divisors of $1$ are those natural numbers $d$ such that $1 = qd$ for some natural number $q$. Since $q \ge 1$ and $d \ge 1$, their product can only be $1$ if both $q = 1$ and $d = 1$. Therefore, $d=1$ is the only positive divisor of $n=1$. Thus, at least one natural number with exactly one positive divisor exists.
            \item (Uniqueness) Suppose $m$ is a natural number that has exactly one positive divisor.
            
            Since the number $1$ divides every natural number, $1$ is a positive divisor of $m$.

            Also, every natural number divides itself. Thus, $m$ is a positive divisor of $m$.

            Since we assumed $m$ has exactly one positive divisor, all divisors of $m$ must be the same number. So the divisor $1$ and the divisor $m$ must be identical: $m=1$.

            Therefore, $m = n = 1$. This proves that the natural number $n$ is unique.
        \end{itemize}
    \end{proof}
    \end{enumerate}
\end{solutions*}

\begin{exercise}
Let $p(x,y)$ be the statement `$x + y$ is even'.
\begin{itemize}
\item Prove that $\forall x \in \mathbb{Z},\, \exists y \in \mathbb{Z},\, p(x,y)$ is true.
\item Prove that $\exists y \in \mathbb{Z},\, \forall x \in \mathbb{Z},\, p(x,y)$ is false.
\end{itemize}
\end{exercise}

\begin{solutions*}
\begin{itemize}
    \item \textbf{Prove that $\forall x \in \mathbb{Z},\, \exists y \in \mathbb{Z},\, p(x,y)$ is true.}

    \begin{proof}
        Let $x$ be an arbitrary integer. Choose $y = x$, then $x + y = 2x$ is even.

        Because we found a valid integer $y$ for arbitrary $x$, the statement is true. More generally,
        \begin{itemize}
            \item If $x$ is even, then $x$ can be represented as $x = 2k$ for some integer $k$. In this case, let $y = 2\ell$ is even and $x + y = 2(k + \ell)$ is even.
            \item If $x$ is odd, then $x$ can be represented as $x = 2k+1$ for some integer $k$. In this case, let $y = 2\ell + 1$ is odd and $x + y = 2(k + \ell + 1)$ is even.
        \end{itemize}
        In summary, regardless of whether $x$ is even or odd, there exists a $y$ such that $x + y$ is even. Therefore, $\forall x \in \mathbb{Z},\, \exists y \in \mathbb{Z},\, p(x,y)$ is true.
    \end{proof}

    \item \textbf{Prove that $\exists y \in \mathbb{Z},\, \forall x \in \mathbb{Z},\, p(x,y)$ is false.}

    \begin{proof}
        To prove this statement is false, the standard mathematical approach is to prove that its logical negation is true.
        
        The negation of $\exists y \in \mathbb{Z},\, \forall x \in \mathbb{Z},\, p(x,y)$ is:
        \[\forall y \in \mathbb{Z},\, \exists x \in \mathbb{Z},\, \neg p(x,y)\]
        Let $y$ be an arbitrary integer. Choose $x = y+1$, then $x + y = 2y+1$ is odd.

        Therefore, for any integer $y$, there exists an integer $x$ that makes $x+y$ odd. Because the negation of the original statement is true, the original statement must be false. 
    \end{proof}
\end{itemize}
\end{solutions*}
